\chapter{Continuous Batching（持续批处理）动态调度仿真}

\section{Lab 12a：显存调度与连续批处理 (Continuous Batching)}

\section{[施工] 为什么静态批处理（Static Batching）在 LLM 中失效？}

在传统的计算机视觉或文本分类任务中，模型输入输出维度是固定的，\textbf{静态批处理（Static Batching）}非常有效。

但在大语言模型生成场景中，每个请求的：
\begin{enumerate}
  \item \textbf{输入 Prompt 长度极不一致}（有的几十字，有的几千字）。
  \item \textbf{生成长度（Output Tokens）难以事先预知}（有的遇到 \texttt{$<$eos$>$} 10 步就结束，有的需要生成 1000 步）。
\end{enumerate}

\subsection{静态 Batching 的两大致命缺陷：}

\begin{lstlisting}[numbers=none]
Request 1: [Prompt: 10] [=== Gen: 20 ===] -> 结束! (闲置等待...) [PADDING PADDING PADDING]
Request 2: [Prompt: 30] [=================== Gen: 100 ===================] -> 最长请求!
Request 3: [Prompt: 5 ] [= Gen: 10 =] -> 结束! (闲置等待...) [PADDING PADDING PADDING PADDING]
\end{lstlisting}

\begin{enumerate}
  \item \textbf{显存与算力气泡（Bubble）}：短请求早早结束，但由于必须作为一个 Batch 统一返回，短请求必须不断 padding 空 token 陪跑，浪费大量算力与带宽。
  \item \textbf{首字延迟排队灾难}：新到达的请求必须等待当前 Batch 中\textbf{最慢的那一个请求完全结束}，才能开始下一批，导致极高的排队排班延迟。
\end{enumerate}

\vspace{0.8em}\hrule\vspace{0.8em}

\section{[加速] 连续批处理 (Continuous Batching / In-flight Batching)}

由 OSDI'22 论文 \textit{Orca} 提出，并在 \textbf{vLLM} 和 \textbf{TensorRT-LLM} 中发扬光大：

\begin{noteBox}[核心导读与背景]
\textbf{核心思想}：调度粒度从「请求级（Request-level）」细化为「迭代步进级（Iteration-level）」。
\end{noteBox}

\begin{lstlisting}[numbers=none]
Iteration t:   [Req 1: decode step 15] [Req 2: decode step 40] [Req 3: decode step 8]
Iteration t+1: [Req 1: 触发 EOS! 退出]   [Req 2: decode step 41] [Req 3: decode step 9]
                ⬇ (槽位立即释放，新请求无缝插入！)
Iteration t+2: [Req 4: 快速 Prefill!]   [Req 2: decode step 42] [Req 3: decode step 10]
\end{lstlisting}

\begin{itemize}
  \item \textbf{立即退场}：任何请求一旦生成结束标记（EOS）或达到最大长度，当步立刻释放槽位与 KV Cache。
  \item \textbf{即时插入}：等待队列中的新请求无需等待其他长请求完成，在下一步前向迭代中直接加入当前 Batch 进行 Prefill。
  \item \textbf{吞吐飞跃}：消除了绝大部分 Padding 气泡，吞吐量显著提升（Orca、vLLM 等论文中常见数倍提升，具体取决于负载）。本仿真会同时打印"Padding 无效计算占比"和"槽位占用率"：前者对连续批处理按构造恒为 0，后者才是两种方案可比的指标。
\end{itemize}

\vspace{0.8em}\hrule\vspace{0.8em}

\section{[文档] 内存革命：PagedAttention 原理}

即使有了 Continuous Batching，连续的内存分配依然会造成严重浪费：
\begin{itemize}
  \item 在传统显存管理中，框架必须为请求预先分配一段\textbf{连续的固定最大长度显存}（如预留 4096 tokens）。
  \item 但大多数请求只生成了 100 个 token。vLLM 论文测到 60\%\textasciitilde{}80\% 的 KV 显存被浪费——这\textbf{不只是内部碎片}，还包括按 max\_len 预留的闲置和外部碎片，三者要分开看（第 12 篇 §3）。
\end{itemize}

\textbf{vLLM 的解决方案}：
借鉴操作系统的\textbf{虚拟内存分页机制（Virtual Memory Paging）}：
\begin{itemize}
  \item 将 KV Cache 切分成固定大小的物理块（Physical Blocks，例如每块 16 个 tokens）。
  \item 物理显存块\textbf{无需连续}，通过一个逻辑块表（Block Table）映射。
  \item 只有当生成了新的 16 个 token 时，才向物理池申请分配 1 个新块。
  \item \textbf{成效}：分页把「最后一个块没写满」以外的浪费基本消掉；单序列的上界是 $15/S$，序列长度 $S \gtrsim 400$ token 时才低于 4\%（短对话达不到，见第 12 篇 §3）；实际收益是同一张卡能承载成倍增加的并发请求数。
\end{itemize}

\vspace{0.8em}\hrule\vspace{0.8em}

\section{[执行] 运行仿真实验}

运行 \textbf{\texttt{continuous\_batching\_sim.py}}：

\begin{lstlisting}[language=bash]
python3 lab12a_continuous_batching/continuous_batching_sim.py
\end{lstlisting}

对比静态批处理与连续批处理在处理相同真实负载时的总吞吐量（Throughput）、平均等待延迟（Latency）与算力浪费（Bubble Ratio）。

\vspace{0.8em}\hrule\vspace{0.8em}

\section{[链接] 继续深入}

\begin{itemize}
  \item \textbf{为什么"每一步耗时视为常数"这个假设成立？} 本仿真假设无论批里有几个请求，一步的耗时都一样。这来自第 8 篇的推导：Decode 线性层的算术强度 $\approx$ batch size，batch 远小于拐点时，一步的耗时由"读一遍权重"决定，与 batch 几乎无关。见 \textbf{第 8 篇 §4} 与 \texttt{lab08\_memory\_and\_roofline/measure\_gpu\_roofline.py} 实验 4。
  \item \textbf{本仿真没有模拟的}：Prefill 的开销、显存上限、长 Prefill 对 Decode 的干扰。这些在 \textbf{第 12 篇} 与 \textbf{Lab 12b}（分页 KV Cache、前缀缓存、Chunked Prefill）中补齐。
\end{itemize}



\section{实验核心源码精选与解析}

本实验配套有完整可运行的工业级验证代码，重点实现细节如下：


\subsection{源码实现：\texttt{continuous\_batching\_sim.py}}

\lstinputlisting[language=Python, caption={\texttt{continuous\_batching\_sim.py}}]{code/lab12a_continuous_batching_continuous_batching_sim.py}

